\documentclass[12pt]{article}
\usepackage[utf8]{inputenc}
\usepackage{hyperref}
\usepackage[margin=1in]{geometry}

\begin{document}

\begin{center}
    \Large\textbf{Base Paper Reference: B92 Protocol}
\end{center}
\vspace{0.5cm}

\noindent\textbf{Title:} Quantum cryptography using any two nonorthogonal states \\
\textbf{Author:} Charles H. Bennett (IBM Research Division) \\
\textbf{Journal:} Physical Review Letters, Vol. 68, Issue 21, pp. 3121--3124 \\
\textbf{Date:} May 25, 1992 \\
\textbf{DOI:} \href{https://doi.org/10.1103/PhysRevLett.68.3121}{10.1103/PhysRevLett.68.3121}

\vspace{0.8cm}
\noindent\Large\textbf{Abstract}
\vspace{0.3cm}\normalsize

\noindent Quantum key distribution (QKD) is typically explained using four quantum states (e.g., the BB84 protocol). In this seminal paper, Charles H. Bennett demonstrates that in principle, \textbf{any two nonorthogonal quantum states} are sufficient for securely distributing a random cryptographic key between two parties. 

\vspace{0.3cm}
\noindent This paper introduces the protocol commonly known today as \textbf{B92}. It establishes the fundamental security proof that eavesdropping on nonorthogonal states necessarily introduces detectable disturbances (Quantum Bit Error Rate), which can be quantified and eliminated through classical post-processing (Information Reconciliation and Privacy Amplification).

\vspace{1cm}
\hrule
\vspace{0.3cm}
\noindent \textit{Note: Due to publisher paywalls (American Physical Society) on this 1992 paper, please use the DOI link above while connected to the university VPN to download the full-text PDF via the institutional library subscription.}

\end{document}
